\section{Theoretical Foundations for Ambulatory Seizure Monitoring}

\subsection{The Ambulatory Monitoring Paradigm: Requirements and Challenges}

\cite{beniczkyAutomatedSeizureDetection2021a,chenSeizuresDetectionUsing2022a}
found that translation of seizure detection from the controlled clinical environment to
an ambulatory setting redifines approaches and performance requirements.
In clinical environments maximum seinsitivity is prioritized over FAR,
when engineering for ambulatory use a balance of both metrics is needed.
Failing to detect a seizure (low sensitivity) poses a safety risk, While
excessive false alarms (high FAR) leads to desensitizing and "alarm fatigue"
and the device being ignored or not used at all.
Users reported that devices should achieve high sensitivity (≥90\%) 
with low false alarm rates — acceptable FARs were 1-2/week for patients with ≥1
seizure/week and 1-2/month for patients with <1 seizure/week
\parencite{sivathambooPreferencesUserExperiences2022}.
While more complex models can find more intricate patterns they also introduce
the "black-box" problem wherein the model lacks Transparency.
Therefore the design should incorporate explanations for decisions to 
ensure users trust and regulatory compliance.
 
\subsection{Physiological Biosignals for Wearable Sensing}

Wearable devices circumvent the impracticality of EEG 
by capturing peripheral biosignals that serve as proxies 
for central nervous system activity during seizures.
These proxies will inherently relay incomplete and 
noisy signals so a combination of multiple sensors
is required. The acceleration/attitude sensor excels at detecting 
muscle convulsions but is noisy in daily use, as it captures all movement. 
The electrodermal activity (EDA) sensor detects autonomic arousal by 
measuring the skin conductance associated with sweat gland activity.

\subsection{Seizure Manifestations and Corresponding Biosignals}

Wearable biosignals act as peripheral proxies for cerebral activity during seizures. 
Different seizure types produce distinct physiological signatures across these signals, 
informing detection approaches. This relationship centers on two primary manifestations: 
motor activity and autonomic activation.

Convulsive seizures, such as generalized tonic-clonic events, 
produce clear motor signatures captured by accelerometry (ACC), 
which measures movement and posture changes characteristic of ictal events 
\parencite{wangEpilepticSeizureDetection2025, wuNovelSeizureDetection2024}.

Simultaneously, seizures drive autonomic responses recorded through complementary biosignals. 
Electrodermal activity (EDA) measures sympathetic nervous system arousal via skin conductance changes 
\parencite{jainCompressedSensingBased2017, mironAutonomicBiosignalsSeizure2025}. 
Cardiac activity provides particularly strong autonomic proxies: electrocardiography (ECG) 
captures heart rate and rhythm alterations 
\parencite{lealViabilityECGFeatures2017, ghaderiAdvancesMachineLearning2025}, 
while heart rate variability (HRV) quantifies autonomic balance dynamics 
\parencite{masonHeartRateVariability2024, paveiEarlySeizureDetection2017}.

In wearable applications, photoplethysmography (PPG) serves 
as a practical alternative for cardiac monitoring,
deriving pulse rate and variability though with greater motion sensitivity than ECG 
\parencite{seizuresDetectionUsingMultimodal, sethFeasibilityCardiacBasedSeizure2023}.

Consequently, detection strategies combine these biosignals according to target seizure types, 
with multimodal fusion providing the most robust approach for ambulatory monitoring.

\subsection{Evaluation Frameworks: From Retrospective Analysis to Prospective Validity}
